\section{Introduction}
\label{introduction}

Developing agents that can perform complex control tasks from high dimensional observations such as pixels has been possible by combining the expressive power of deep neural networks with the long-term credit assignment power of reinforcement learning algorithms. Notable successes include learning to play a diverse set of video games from raw pixels \cite{mnih2015human}, continuous control tasks such as controlling a simulated car from a dashboard camera \cite{lillicrap2015continuous} and subsequent algorithmic developments and applications to agents that successfully navigate mazes and solve complex tasks from first-person camera observations \cite{jaderberg2016reinforcement, espeholt2018impala, jaderberg2019human}; and robots that successfully grasp objects in the real world \cite{kalashnikov2018qt}. 


\begin{figure}[!ht]
\vskip 0.2in
\begin{center}
\centerline{\includegraphics[width=7cm]{plots/curl_momentum_diagram.png}}

\caption{\textbf{C}ontrastive \textbf{U}nsupervised Representations for \textbf{R}einforcement \textbf{L}earning (CURL) combines instance contrastive learning and reinforcement learning. CURL trains a visual representation encoder by ensuring that the embeddings of data-augmented versions $o_q$ and $o_k$ of observation $o$ match using a contrastive loss. The \emph{query} observations $o_q$ are treated as the anchor while the \emph{key} observations $o_k$ contain the positive and negatives, all constructed from the minibatch sampled for the RL update. The keys are encoded with a momentum averaged version of the query encoder. The RL policy and (or) value function are built on top of the query encoder which is jointly trained with the contrastive and reinforcement learning objectives. CURL is a generic framework that can be plugged into any RL algorithm that relies on learning representations from high dimensional images.}



\label{momentum_encoder}
\end{center}
\vskip -0.2in
\end{figure}


However, it has been empirically observed that reinforcement learning from high dimensional observations such as raw pixels is sample-inefficient \cite{lake2017building, kaiser2019model}. Moreover, it is widely accepted that learning policies from physical state based features is significantly more sample-efficient than learning from pixels \cite{tassa2018deepmind}. In principle, if the state information is present in the pixel data, then we should be able to learn representations that extract the relevant state information. For this reason, it may be possible to learn from pixels as fast as from state given the right representation.

From a practical standpoint, although high rendering speeds in simulated environments enable RL agents to solve complex tasks within reasonable wall clock time, learning in the real world means that agents are bound to work within the limitations of physics. \citet{kalashnikov2018qt} needed a farm of robotic arms that collected large scale robot interaction data over several months to develop their robot grasp value functions and policies. The data-efficiency of the whole pipeline thus has significant room for improvement. Similarly, in simulated worlds which are limited by rendering speeds in the absence of GPU accelerators, data efficiency is extremely crucial to have a fast experimental turnover and iteration. Therefore, improving the sample efficiency of reinforcement learning (RL) methods that operate from high dimensional observations is of paramount importance to RL research both in simulation and the real world and allows for faster progress towards the broader goal of developing intelligent autonomous agents.

A number of approaches have been proposed in the literature to address the sample inefficiency of deep RL algorithms. Broadly, they can be classified into two streams of research, though not mutually exclusive: (i) Auxiliary tasks on the agent's sensory observations; (ii) World models that predict the future. While the former class of methods use auxiliary self-supervision tasks to accelerate the learning progress of model-free RL methods \cite{jaderberg2016reinforcement, mirowski2016learning}, the latter class of methods build explicit predictive models of the world and use those models to plan through or collect fictitious rollouts for model-free methods to learn from \cite{sutton1990integrated, ha2018world, kaiser2019model, schrittwieser2019mastering}.


Our work falls into the first class of models, which use auxiliary tasks to improve sample efficiency. Our hypothesis is simple: {\it If an agent learns a useful semantic representation from high dimensional observations, control algorithms built on top of those representations should be significantly more data-efficient.} Self-supervised representation learning has seen dramatic progress in the last couple of years with huge advances in masked language modeling \cite{devlin2018bert} and contrastive learning \cite{henaff2019data, he2019momentum, chen2020simclr} for language and vision respectively. The representations uncovered by these objectives improve the performance of any supervised learning system especially in scenarios where the amount of labeled data available for the downstream task is really low. 


We take inspiration from the contrastive pre-training successes in computer vision. However, there are a couple of key differences: (i) There is no giant unlabeled dataset of millions of images available beforehand - the dataset is collected online from the agent's interactions and changes dynamically with the agent's experience; (ii) The agent has to perform unsupervised and reinforcement learning simultaneously as opposed to fine-tuning a pre-trained network for a specific downstream task. These two differences introduce a different challenge: How can we use contrastive learning for improving agents that can learn to control effectively and efficiently from online interactions? 


To address this challenge, we propose CURL - \textbf{C}ontrastive \textbf{U}unsupervised Representations for \textbf{R}einforcement \textbf{L}earning. CURL uses a form of contrastive learning that maximizes agreement between augmented versions of the same observation, where each observation is a stack of temporally sequential frames. We show that CURL significantly improves sample-efficiency over prior pixel-based methods by performing contrastive learning simultaneously with an off-policy RL algorithm. CURL coupled with the Soft-Actor-Critic (SAC) \cite{haarnoja2018soft} results in {\textbf{1.9x}} median higher performance over Dreamer, a prior state-of-the-art algorithm on DMControl environments, benchmarked at {\textbf{100k}} {\it environment steps} and {\it matches the performance of state-based SAC} on the majority of 16 environments tested, a {\textbf{first}} for pixel-based methods. In the Atari setting benchmarked at 100k {\it interaction steps}, we show that CURL coupled with a data-efficient version of Rainbow DQN \cite{van2019use} results in {\textbf{1.2x}} median higher performance over prior methods such as SimPLe~\cite{kaiser2019model}, improving upon Efficient Rainbow~\cite{van2019use} on {\it 19 out of 26} Atari games, {\it surpassing human efficiency} on two games. 

While contrastive learning in aid of model-free RL has been studied in the past by \citet{oord2018representation} using Contrastive Predictive Coding (CPC), the results were mixed with marginal gains in a few DMLab \cite{espeholt2018impala} environments. CURL is the first model to show substantial data-efficiency gains from using a contrastive self-supervised learning objective for model-free RL agents across a multitude of pixel based continuous and discrete control tasks in DMControl and Atari.


We prioritize designing a simple and easily reproducible pipeline. While the promise of auxiliary tasks and learning world models for RL agents has been demonstrated in prior work, there's an added layer of complexity when introducing components like modeling the future in a latent space \cite{oord2018representation, ha2018world}. CURL is designed to add minimal overhead in terms of architecture and model learning. The contrastive learning objective in CURL operates with the same latent space and architecture typically used for model-free RL and seamlessly integrates with the training pipeline without the need to introduce multiple additional hyperparameters. 

Our paper makes the following {\textbf{key contributions}}: We present CURL, a simple framework that integrates contrastive learning with model-free RL with minimal changes to the architecture and training pipeline. Using 16 complex control tasks from the DeepMind control (DMControl) suite and 26 Atari games, we empirically show that contrastive learning combined with model-free RL outperforms the prior state-of-the-art by 1.9x on DMControl and 1.2x on Atari compared across leading prior pixel-based methods. CURL is also the first algorithm {\it across both model-based and model-free methods} that operates purely from pixels, and nearly matches the performance and sample-efficiency of a SAC algorithm trained from the state based features on the DMControl suite. Finally, our design is simple and does not require any custom architectural choices or hyperparameters which is crucial for reproducible end-to-end training.  Through these strong empirical results, we demonstrate that a contrastive objective is the preferred self-supervised auxiliary task for achieving sample-efficiency compared to reconstruction based methods, and enables {\it model-free methods to outperform state-of-the-art model-based methods in terms of data-efficiency}.
